\setcounter{footnote}{6}

\printedpage{9}
% head: A. Grothendieck : Dérivateurs, Ch. XVIII, Édition des Universités de Montpellier II et Paris VII

\noindent\textit{Moreover, if c is satisfied, we have this: for $s\in S$, $t\in T$, let us consider}
\[
(s,t)\ \in\ (S\times T)\ \subset\ (S\times T)^\wedge\ \simeq\ S^\wedge\boxtimes T^\wedge\ ,
\]
\textit{which is also written out as}
\[
\left\{
\begin{aligned}
&\rlap{$\displaystyle (s,t)\ =\ \pr_1^*(s)\times\pr_2^*(t)\ \simeq\ s\overset{S^\wedge,T^\wedge}{\boxtimes}t\in\Ob\,(S\times T)^\wedge\ ,$}\\
&\textit{i.e.}\ (s,t)(s',t')\ =\ s(s')\times t(t') &&=\ \Hom_S(s',s)\times\Hom_T(t',t)\\
& &&=\ \Hom_{S\times T}((s',t'),(s,t))\ ,
\end{aligned}
\right.
\tag*{(4.3.1.4)}
\]
\[
\begin{tikzcd}
S\times T \arrow[r, "\mathrm{pr}_2"] \arrow[d, "\mathrm{pr}_1"] & T\\
S\ ,
\end{tikzcd}
\]
\msp{page 66}\textit{we have}
% note: the label (4.3.1.4) is printed for this display as well as for the preceding one
\[
s\overset{\mathcal M,\mathcal N}{\boxtimes}t\ \simeq\ \varphi(s,t)\ ,
\tag*{(4.3.1.4)}
\]
\textit{and $\varphi(S\times T)=\{s\boxtimes t\mid s\in S,\ t\in T\}$ is a strictly generating subset of $\mathcal M\boxtimes\mathcal N$ (}\footnote{\textbf{N.B.} Of course we have $(s,t)=s\overset{S^\wedge,T^\wedge}{\boxtimes}t$.}\textit{).}

\noindent\textbf{N.B.} I know of no example where the conditions under consideration are not satisfied, but nevertheless now doubt that even the weakest always is, i.e. that $\mathcal M\boxtimes\mathcal N$ is a pseudo-topos as soon as $\mathcal M$ and $\mathcal N$ are.

\noindent\textsc{Proof of} 4.3.1. The implication b $\Longrightarrow$ b$'$ is tautological. The implication b$'$ $\Longleftrightarrow$ c follows from the

% marginal vertical bar printed beside Lemme 4.3.2 (from "Lemme 4.3.2." to "soit représentable.")
\noindent\textbf{Lemma 4.3.2.} \textit{Let $\mathcal E$, $X$ be $\mathfrak U$-categories with $X$ small, $\mathcal E$ stable under small $\varprojlim$ ,}
\[
\alpha\ :\ \mathcal E\longrightarrow X^\wedge
\]
\textit{a functor. For $\alpha$ to admit a left adjoint $\varphi$, it is necessary and sufficient that for every $x\in\Ob X$, the functor}
\[
\left\{
\begin{array}{ccl}
\xi & \longmapsto & \Hom_{X^\wedge}(x,\alpha(\xi))\ =\ \alpha(\xi)(x)\\
\mathcal E^\circ & \longrightarrow & \mathrm{Ens}
\end{array}
\right.
\]
\textit{be representable.}

\noindent This is obviously necessary. Sufficiency follows from the fact that $X\subset X^\wedge$ is strict%
% foot: Version 17 décembre 2025
\printedpage{10}%
% head: A. Grothendieck : Dérivateurs, Ch. XVIII, Édition des Universités de Montpellier II et Paris VII
\msp{page 67}%
ly generating in $X^\wedge$, i.e. that every object $F$ of $X^\wedge$ is a direct limit $F=\varinjlim\limits_I F_i$, the $F_i$ in $X$ (we take $I=X/F$), whence
\[
\begin{aligned}
\Hom(F,\alpha(\xi))\ &\simeq\ \varprojlim_I\ \underbrace{\Hom(F_i,\alpha(\xi))}_{\Hom(\rho(F_i),\xi)}\\
&\simeq\ \Hom(\varinjlim_I\,(\rho(F_i),\xi))\ ,
\end{aligned}
\]
so the functor $\xi\longmapsto\Hom(F,\alpha(\xi))$ is represented by $\rho(F)\overset{\text{def}}{=}\varinjlim\limits_I\rho(F_i)$.

\noindent Let us prove c $\Longrightarrow$ a. For the existence of small $\varinjlim$ in $\mathcal M\boxtimes\mathcal N$, and the existence of a strictly generating subset, we invoke the general lemma:

% marginal vertical bar printed beside Lemme 4.3.3 (from "Lemme 4.3.3." to "comme")
\noindent\textbf{Lemma 4.3.3.} \textit{Let $\alpha:\mathcal E\hookrightarrow\mathcal F$ be a fully faithful functor admitting a left adjoint $\rho$.}
\begin{enumerate}[label={}]
\item[\textcircled{a}] \textit{If the $\varinjlim$ of a given type $I$ exist in $\mathcal F$, they exist in $\mathcal E$, and are computed as}
\[
\varinjlim_I{}_{\mathcal E}\,\xi_i\ =\ \rho(\varinjlim_I{}_{\mathcal F}\,\alpha(\xi_i))\ .
\tag*{(4.3.3.1)}
\]
\item[\textcircled{b}] \textit{If $U$ is a small strictly generating subset of $\mathcal F$, $\rho(U)$ is a small strictly generating subset in $\mathcal F$.}
\end{enumerate}

\noindent Indeed, if $\xi$ denotes the right-hand side \texttt{[of equality 4.3.3.1]}, and $\eta$ a variable object of $\mathcal E$, we will have
\[
\begin{aligned}
\Hom_{\mathcal E}(\xi,\eta)\ &\overset{\text{adj.}}{\simeq}\ \Hom_{\mathcal F}(\varinjlim_I\alpha(\xi_i),\alpha(\eta))\\
&\simeq\ \varprojlim_I\quad\underbrace{\Hom_{\mathcal F}(\alpha(\xi_i),\alpha(\eta))}_{\simeq\,\Hom_{\mathcal E}(\xi_i,\eta)\ \ \text{since $\alpha$ fully faithful}}\quad .
\end{aligned}
\]
\msp{page 68}
\[
\simeq\ \ \varprojlim_I\Hom_{\mathcal E}(\xi_i,\eta)\ ,
\]
which proves that $\xi$ is a direct limit of the $\xi_i$ in $\mathcal E$ (\footnote{To prove b), we note that the full faithfulness of $\alpha$ implies that $\rho$ is essentially surjective; as $\rho$ commutes with small $\varinjlim$ and every object of $\mathcal F$ is of the form $\varprojlim\limits_I x_i$, the $x_i$ in $U$, we will have $\rho(x)=\varinjlim\limits_I\rho(x_i)$, the $\rho(x_i)$ in $\rho(U)$, we win.}).

\noindent This proves that (assuming c)) $\mathcal M\boxtimes\mathcal N$ is stable under small $\varinjlim$ , so it remains to prove that $\mathcal M\boxtimes\mathcal N$ admits a small strictly generating family. This will follow from b) in the lemma, since $S\times T$ strictly generates $(S\times T)^\wedge$. This proves at the same time the last assertion of the proposition, subject to proving formula (4.1.3.4). We must therefore prove an isomorphism functorial in $F\in\Ob\mathcal M\boxtimes\mathcal N$, % [sic] (4.1.3.4) -> (4.3.1.4), the formula s ⊠ t ≃ φ(s,t) to be proved is labelled (4.3.1.4) on p. 9
\[
\Hom(\rho(s,t),F)\ \simeq\ F(s,t)\ ,
\]
but by adjunction we have
\[
\Hom(\rho(s,t),F)\ \simeq\ \Hom((s,t),\alpha(F))\ =\ \alpha(F)(s,t)\ =\ F(s,t)\ ,
\]
OK.

% foot: Version 17 décembre 2025
\printedpage{11}
% head: A. Grothendieck : Dérivateurs, Ch. XVIII, Édition des Universités de Montpellier II et Paris VII

\noindent\msp{page 69}\textbf{Proposition 4.3.4.} \textit{Let $\mathcal M$ be a $\mathfrak U$-category. The following conditions are equivalent.}
\begin{enumerate}[label={}]
\item[\textit{a)}] \textit{$\mathcal M$ is a $\mathfrak U$-pseudo-topos category, i.e. it is stable under $\mathfrak U$-direct limits, and admits a $\mathfrak U$-small strictly generating subcategory $S$.}
\item[\textit{b)}] \textit{There exist a $\mathfrak U$-small category $S$, and a functor $i:\mathcal M\hookrightarrow S^\wedge$ that is} fully faithful, \textit{having a left adjoint.}
\end{enumerate}

\noindent\textsc{Proof}. We have a) $\Longrightarrow$ b), for it suffices to take for $S$ in b) the one specified in a), the hypothesis on $S$ implies that the canonical functor
\[
\mathcal M\ \longrightarrow\ S^\wedge\ ,\qquad\quad \text{composite of } \mathcal M\hookrightarrow\ \mathcal M^\wedge\ \xrightarrow{\ \text{restriction}\ }S^\wedge
\]
is fully faithful. That it admits a left adjoint follows from lemma 4.3.2. The implication b) $\Longrightarrow$ a) follows from lemma 4.3.3.

\noindent\textbf{Corollary 4.3.5.} \textit{Let $\mathcal M$ be a $\mathfrak U$-category. For $\mathcal M$ to be a $\mathfrak U$-pseudo-topos that is} accessible \textit{(i.e. all of whose objects are accessible), it is necessary and sufficient that there exist $S$ in} Cat \textit{and a functor fully faithful and} accessible $\mathcal M\overset{\alpha}{\hookrightarrow}S^\wedge$\textit{, admitting a left adjoint $\rho$.}

\noindent\msp{page 70}\textsc{Proof}. It suffices to note that if $S$, $\alpha$ are as in 4.3.4 b), then $\alpha$ is accessible if and only if the $\rho(S)$ are made up of accessible objects in $\mathcal M$; as this subset is strictly generating, we know that this implies that \textit{every} object of $\mathcal M$ is accessible. (Cf. SGA 4 I 9.9.)

\noindent I am going to try to give a more complete statement:

\noindent\textbf{Theorem 4.3.6.} \textit{Let $\mathcal M$ be a $\mathfrak U$-category. Equivalent conditions:}
\begin{enumerate}[label={}]
\item[\textit{a)}] \textit{$\mathcal M$ stable under small $\varinjlim$ , and it admits a small full subcategory $S$, strictly generating, and consisting of accessible objects of $\mathcal M$.}
\item[\textit{a$'$)}] \textit{$\mathcal M$ is a weakly accessible $\mathfrak U$-pseudo-topos category, i.e. satisfies a), with in addition} every \textit{element of $\mathcal M$ (not only those of $S$) accessible in $\mathcal M$.}
\item[\textit{b)}] \textit{There exist a small category $S$, and a functor $\alpha:\mathcal M\hookrightarrow S^\wedge$ that is} fully faithful and accessible, \textit{such that $\alpha$ admits a left adjoint.}
\item[\textit{c)}] \msp{page 71}\textit{$\mathcal M$ is stable under small $\varinjlim$ , is accessible, and admits a cardinal filtration} $(\Filt^\pi\mathcal M)_{\pi\geqslant\pi_0}$.
\item[\textit{d)}] \textit{There exist a cardinal $\pi_1$, a small category $\mathcal M_{\pi_1}$, stable under direct limits of cardinality $\leqslant\pi_1$, such that we have an equivalence of categories}
\[
\mathcal M\ \xleftarrow{\ \sim\ }\ \Ind_{\pi_1}\mathcal M_{\pi_1}\ ,
\tag*{(4.3.6.1)}
\]
\textit{where the right-hand side denotes the full subcategory of} $\Ind(\mathcal M_{\pi_1})$ \textit{consisting of the ind-objects $(X_i)_{i\in I}$ of $\mathcal M_{\pi_1}$ such that $I$ is large relative to $\pi_1$.}
\end{enumerate}

\noindent\textsc{Proof of} 4.3.6 (\footnote{For a complete proof (the one given here is botched from a certain point on), cf. §4.6 (p. 139-148).}). We have already seen a$'$ $\Longrightarrow$ a $\Longrightarrow$ b $\Longrightarrow$ a$'$, i.e. the equivalence of a, a$'$, b, let PTA be this condition. We are going to prove the implications
\[
\mathrm{PTA}\ \Longrightarrow\ \mathrm{c}\ \Longrightarrow\ \mathrm{d}\ \Longrightarrow\ \mathrm{PTA}\ ,
\]
% foot: Version 17 décembre 2025
\printedpage{12}%
% head: A. Grothendieck : Dérivateurs, Ch. XVIII, Édition des Universités de Montpellier II et Paris VII
which will establish the theorem.

\noindent\textit{Let us prove PTA $\Longrightarrow$ c.} Let $S$ be a small strictly generating full subcategory. Thus, \msp{page 72}for every $x$ in $\mathcal M$ we have
\[
x\ =\ \varinjlim_{s\in S/x}s
\tag*{(4.3.6.2)}
\]
(cf. SGA 4 I 7.2 (i)). Let $\pi_0=\card\Fl(S)$, and for every cardinal $\pi\geqslant\pi_0$, let
\[
\begin{array}{lcl}
\Filt^\pi\mathcal M & = & \text{full subcategory of $\mathcal M$ consisting of objects $x$ such that}\\
 & & \card\,(\Fl(S/x))\leqslant\pi\ .
\end{array}
\tag*{(4.3.6.3)}
\]
I want to prove that we thus get a cardinal filtration on $\mathcal M$, i.e. that conditions a) b) c) of SGA 4 I 9.12 are satisfied. This is proved in loc. cit. 9.13 under slightly different hypotheses, the condition of accessibility of $\mathcal M$ being replaced by the one (much less pleasant!) that the strict epimorphic families of arrows in $\mathcal M$ are strict universal epimorphic families. This condition is used only in the proof of condition b), namely the stability of $\Filt^\pi\mathcal M$ under $\varinjlim\limits_I$ of cardinality $\leqslant\pi$ (without, moreover, having to assume that $I$ is filtering, nor a fortiori large relative to $\pi_0$). It must moreover be said that the definition of $\Filt^\pi\mathcal M$ given in loc. cit. is not quite the same as (4.3.6.3), \msp{page 73}but it is
% the printed line break "car-/dinal" inside the display is kept as printed
\[
\begin{array}{lcl}
\Filt'^{\pi}\mathcal M & = & \text{full subcategory of $\mathcal M$ consisting of objects $x$ such that there}\\
 & & \text{exists a strict epimorphic family $x_i\xrightarrow{\alpha_i}x$ of car-}\\
 & & \text{dinality $\leqslant\pi$, with $x_i\in S\ \forall\ i\in I$.}
\end{array}
\tag*{(4.3.6.3$'$)}
\]
Thus we have
\[
\Filt^\pi\mathcal M\ \subset\ \Filt'^{\pi}\mathcal M\ ,
\]
but in loc. cit. it is proved that for $x$ in $\Filt'^{\pi}\mathcal M$, we have
\[
\card\Fl(S/x)\ \leqslant\ \pi^{\pi_0}\ ,
\]
whence
\[
\Filt^\pi\mathcal M\ \subset\ \Filt'^{\pi}\mathcal M\ \subset\ \Filt^{\pi^{\pi_0}}\mathcal M\ ,
\]
so if $\pi=\pi^{\pi_0}$, e.g. $\pi=2^c$ with $c\geqslant\pi_0$, then we get
\[
\Filt^\pi\mathcal M\ =\ \Filt'^{\pi}\mathcal M\ .
\]

\noindent So the two filtrations look $\pm$ equivalent. But the proof given of the stability of $\Filt'^{\pi}$ really seems to make essential use of the universality condition on strict epimorphic families (hence effective epimorphic, with the hypotheses made implying the existence of small $\varprojlim$ ). So I must, willy-nilly, do the work over again! % [sic] on l'air -> ont l'air, verb avoir (3rd person plural) intended

% foot: Version 17 décembre 2025
\printedpage{13}
% head: A. Grothendieck : Dérivateurs, Ch. XVIII, Édition des Universités de Montpellier II et Paris VII

\noindent\msp{page 74}Condition a): $\Filt^\pi$ equivalent to a small category. This is immediate, since by (4.3.6.2), $\Filt^\pi$ consists of $\varprojlim\limits_I$ in $\mathcal M$ of objects of $S$ (a small category), with $I$ such that $\card\Fl I\leqslant\pi$.

\noindent Condition b): $\Filt^\pi$ stable under $\varprojlim\limits_I$ , when $\card\Fl I\leqslant\pi$ (\footnote{when $I$ filtering and large relative to $\pi_0$ -- I prove it only in this case, alas!}). So let $(x_i)_{i\in I}$, where
\[
\begin{array}{ccc}
I & \longrightarrow & \mathcal M\\[1ex]
i & \longmapsto & x_i
\end{array}
\]
a direct system in $\mathcal M$, with $x_i$ in $\Filt^\pi$. Let us consider the fiber product
\[
\begin{tikzcd}
I \arrow[d, "\varphi"'] & S/I = J \arrow[l, "\text{cofibered}"'] \arrow[d] \arrow[dr, dotted, "\psi"] & & \\
\mathcal M & S/\mathcal M \arrow[l, "\text{target}"', "\text{Cat-cofibered}"] \arrow[r, "\text{source}"] & S \arrow[r, hook] & \mathcal M\ ,
\end{tikzcd}
\tag*{(4.3.6.4)}
\]
where $J$ is Cat-cofibered over $I$, its fiber at $i\in\Ob I$ being isomorphic to $S/x_i$, so $\card\Fl J_i\leqslant\pi$. Since by hypothesis $\card\Fl I\leqslant\pi$, we conclude that $\card\Fl J\leqslant\pi^3=\pi$. (\textbf{N.B.} The set of arrows $v:j_0\longrightarrow j_1$ of $J$ over an arrow $u:i_0\longrightarrow i_1$ of $I$ is ``contained'' in the product set $\Ob J_{i_0}\times\Fl J_{i_1}$, by the map
\[
v\longmapsto(\operatorname{source}v,v':u_*(j_0)\longrightarrow j_1)\ ,
\]
hence of cardinality bounded above by $\pi^2$.) We \msp{page 75}then have
\[
x\ \overset{\text{def}}{=}\ \varinjlim_I x_i\ =\ \varinjlim_J\ \underbrace{\xi_j}_{=\,\psi(j)}\qquad\quad\mbox{(\footnotemark)}\ ,
\]
\footnotetext{We use the distributivity property of $\varinjlim$ , relative to a Cat-cofibrant functor
\[
\begin{tikzcd}[ampersand replacement=\&]
J \arrow[r] \arrow[dr] \& I\\
\& \mathcal M\ ,
\end{tikzcd}
\]
made explicit above.}%
where this time the limit is a limit of objects of $S$, taken in $\mathcal M$. It follows that $x\in\Filt'^{\pi}$, this is not quite $x\in\Filt^\pi$. (That implies $x\in\Filt^{\pi^{\pi_0}}$, so $x\in\Filt^\pi$ if $\pi=\pi^{\pi_0}$, e.g. $\pi=2^c$ with $c\geqslant\pi_0$.) But this is the moment to use the fact that the objects of $S$ are accessible, and we will take $\pi_0$ such that the objects of $S$ are $\pi_0$-accessible, and this time we will take $I$ filtering and large relative to $\pi_0$ all the same (to restrict ourselves to what is required in condition b) of loc. cit.). We must prove that $\card\,(\Fl(S/x))\leqslant\pi$ (which means $x\in\Filt^\pi$). But since the objects of $S$ are $\pi_0$-accessible and $I$ large relative to $\pi_0$, we see at once that
\[
S/x\ \simeq\ \varinjlim_i S/x_i\ ,
\]
% foot: Version 17 décembre 2025